\chapter{Reference-Adjusted Summarization Evaluation} \label{chap:chap-4}



LLMs have been shown to be effective judges of generations across a range of tasks, from summarization to machine translation. Although they perform well in reference-free settings, evidence suggests that providing a reference for comparison further improves their judgment. However, exact references are rarely available without expensive dataset curation, which requires collecting gold references for each example in the test dataset. In this paper, we propose a new evaluation paradigm to address this gap: reference-adjusted evaluation. This refers to the setting in which a related, but not exact, reference is available. We show that providing related summaries improves LLM judgment by giving them an axis of comparison. This setting enables researchers and developers to improve the evaluations of new test sets without requiring expensive annotations.


% \section{Introduction}
% While human evaluation remains the gold standard for summarization evaluation, automatic evaluation metrics have played an important role in summarization research as they allow for inexpensive and quick iteration for model development. Conventional metrics such as ROUGE~\cite{Lin2004ROUGEAP} or BLEU~\cite{Papineni2002BleuAM} have a number of well-documented issues, such as low correlation with human-judgment~\cite{Liu2008CorrelationROUGE,Fabbri2021SummEval}. Given the promising capabilities of Large Language Models (LLMs), recent research has explored used LLMs for evaluation, a paradigm known as LLM-as-judges~\cite{Zheng2023JudgingLLM,Liu2023GEval}. Research in LLMs-as-judges shows potential for high correlation with human judgment and the ability to generalize to a variety of text generation topics~\cite{Zheng2023JudgingLLM,Liu2023GEval}. The two current paradigms of evaluation metrics are (1) reference-based and (2) reference-free.
% LLMs-as-judges perform best when provided with a reference~\cite{Zheng2023JudgingLLM}. However, collecting gold reference data is expensive, sparking interest in reference-free evaluation metrics that judge summaries solely on the output summary~\cite{Deutsch2022LimitationsReferenceFree,Liu2023GEval}.


% \textbf{We propose a new evaluation paradigm: reference-adjusted evaluation.} This setting refers to the situation in which we have access to similar, but not exact reference summaries. For example, when generating lay summaries of scientific documents, we often have access to abstracts (expert-targeted summaries) but not gold lay summaries. These summaries share the desired content, but not style of the target summaries. Alternatively, say we have access to gold lay summaries, but not one for each test document; these summaries share the target style, but not content. We hypothesize that having related but non-exact summaries will improve LLMs as judges by giving the model an axis on which to refer. This setting will allow researchers and developers to improve their evaluations in new test sets without requiring expensive annotations, leveraging either a small annotation (a few new style examples) or naturally existing data (such as abstracts in the example above). We show that reference-adjusted evaluation can leverage cheaply-acquired machine generated summaries, including extractive summaries, and improve over the reference-free setting. This finding will allow researchers to improve their LLM-as-judge evaluations in situations in which gold reference data is not available or prohibitively expensive.



\section{Introduction}

Automatic evaluation metrics are a cornerstone of summarization research,
enabling rapid and inexpensive iteration during model development. Yet the two
dominant paradigms, reference-based and reference-free evaluation,
each carry a fundamental limitation that constrains their practical utility.

% Reference-based metrics such as ROUGE~\cite{Lin2004ROUGEAP} and
% BLEU~\cite{Papineni2002BleuAM} measure token overlap against a gold
% human-written summary, but gold summaries are expensive to collect: they
% require human annotators to write or validate a reference for every document
% in every new test set. This cost becomes prohibitive when evaluating on new
% domains, new styles, or new tasks where no prior annotations exist. Moreover,
% these metrics exhibit well-documented low correlation with human
% judgment~\cite{Liu2008CorrelationROUGE,Fabbri2021SummEval}, limiting their
% reliability even when references are available.

Given recent success building high-performance Large Language Models (LLMs), researchers have also explored using LLMs to evaluate the generated output of other LLMs, a setting known as LLM-as-Judge. LLM-as-judge evaluation~\cite{Zheng2023JudgingLLM,liu-etal-2023-g}
leverages the reasoning capabilities of LLMs to score generated summaries,
achieving high correlation with human judgment across a range of text
generation tasks. Crucially, LLM judges can operate in a
\textit{reference-free} setting~\cite{Deutsch2022LimitationsReferenceFree},
assessing summary quality without any gold reference. This completely removes the
annotation bottleneck. However, reference-free LLM judges
lack an external anchor for their assessments, and prior work shows that they
perform substantially better when a reference is
provided~\cite{Zheng2023JudgingLLM,Kim2023Prometheus}.

This creates a practical dilemma for researchers and developers. Gold
references yield better evaluations, but collecting them is expensive.
Reference-free evaluation is cheap but weaker. In many real-world settings,
this tradeoff is particularly acute: a researcher evaluating a new
domain-adapted summarization model may have access to \textit{related}
summaries, such as abstracts of scientific articles, or a handful of
human-written examples from a pilot annotation, but not a gold reference
for every test document. Under current paradigms, this approximate information is not leveraged. 




\textbf{We propose reference-adjusted evaluation}, a new paradigm that
occupies the practical middle ground between the two current standards. Rather than requiring exact gold references or discarding all reference information
entirely, reference-adjusted evaluation leverages \textit{related but
non-exact} references: summaries that share either the desired content or the
target style of the ideal output, but not necessarily both. For example, when
evaluating lay summaries of scientific articles, document abstracts are
naturally available. They share the desired informational content but are
written for an expert audience rather than a lay one. Similarly, a small set
of gold lay summaries from a pilot study shares the target style but covers
different documents than the test set. In both cases, the related summary
provides a useful reference point for the LLM judge without requiring a
full annotation effort. 

We introduce the reference-adjusted evaluation paradigm and demonstrate that
it improves LLM-as-judge performance over the reference-free setting across
three open-source models and four evaluation axes, using cheaply acquired
machine-generated summaries as approximate references. We highlight a practical path forward for researchers who need robust evaluation without the cost of full gold annotation. We visualize the trade-offs of the 3 evaluation paradigms in~\Cref{fig:ref-adj-vis}. Reference-based evaluation has high correlation with human judgment, but lacks scalability and ease of deployment, while requiring high annotation cost. Reference-free evaluation is scalable, easy to deploy, and requires no annotation, but suffers from lower correlation with human judgment. Reference-adjusted evaluation bridges this gap by improving easy of deployment, scalability, low to no annotation cost, and higher correlation with human judgment. 
\section{Related Works}
\label{sec:related-work}

The most common evaluation setup for summarization research is to use metrics such as ROUGE~\cite{Lin2004ROUGEAP} or BLEU~\cite{Papineni2002BleuAM} that measure token overlap with a gold reference. These metrics are favored for their interpretability and established use in summarization research, allowing for easy comparison with past work. However, these metrics have a number of well documented issues, such as showing low correlation with human judgment~\cite{Liu2008CorrelationROUGE,fabbri-etal-2021-summeval}. Therefore, these metrics are often combined with other measures, such as BERTScore, which uses a masked language model to measure semantic overlap with a reference summary~\cite{zhang2020bertscoreevaluatingtextgeneration}. Researchers have focused on designing metrics that better correlate with human judgments of quality~\cite{zhang2020bertscoreevaluatingtextgeneration,Fabbri2021SummEval}. However, there is little work in reproducing studies on evaluation metrics, despite our field's reliance on their outcomes. We include a number of traditional metrics in both our reproduction study and as baselines for our LLM-as-judge experiments. 

 Recent work has shown great success using LLMs-as-judges, showing high correlation with human judgments across a variety of tasks and domains~\cite{Zheng2023JudgingLLM}. Researchers have improved LLM-as-Judge performance through methods such as evaluation-specific finetuning~\cite{Kim2023Prometheus} or better prompting~\cite{liu-etal-2023-g}. Although LLMs-as-Judges perform best when provided with a gold reference~\cite{Deutsch2022LimitationsReferenceFree}, there is great interest in the reference-free setting of evaluation metrics as gold references are often unavailable or prohibitively expensive to collect~\cite{liu-etal-2023-g,vasilyev-etal-2020-fill}. We propose a new setting, reference-adjusted evaluation, in which relevant but non-exact references are available.
\input{research-chapters/6-chapter/3-Method}

\section{Experimental Setup}\label{sec:6-chapter-experimental-setup}

We use the SummEval dataset as the basis of our analysis~\cite{fabbri-etal-2021-summeval}. SummEval is based on the CNN/DM dataset~\cite{Hermann2015TeachingMT,Nallapati2016AbstractiveTS}. It includes human annotations for summaries generated by 16 models over 100 source documents, totaling 1600 annotated summaries. The generated summaries are annotated on 4 axes: coherence, consistency, fluency, and relevance.  Each document has 11 reference summaries. We choose this dataset as it is most comprehensive collection of human annotations of model generations and is the standard for measuring correlation with human judgment for summarization evaluation metrics~\cite{fabbri-etal-2021-summeval,liu-etal-2023-g}. We report the Kendall-Tau correlation with the expert annotation judgments.

We follow~\textcite{fabbri-etal-2021-summeval} include a collection of traditional evaluation metrics as baselines. These represent the non-LLM-based metrics. We experiment with the following metrics: ROUGE (-1,-3,-4,-L,-su, -w)~\cite{Lin2004ROUGEAP},  BertScore~\cite{Zhang2019BERTScoreET},  SummQA~\cite{scialom-etal-2019-answers}, BLANC~\cite{vasilyev-etal-2020-fill}, SUPERT~\cite{gao-etal-2020-supert}, METEOR~\cite{banerjee-lavie-2005-meteor}, Lexical Statistics~\cite{fabbri-etal-2021-summeval}, CIDer~\cite{Vedantam2015CIDEr}, CHRF~\cite{popovic-2015-chrf}, and BLEU~\cite{Papineni2002BleuAM}. We use the SummEval toolkit to calculate these metrics~\cite{fabbri-etal-2021-summeval}.  We exclude the following metrics from the original SummEval study: Sentence Mover's Similarity (SMS)~\cite{clark-etal-2019-sentence}, ROUGE-we~\cite{ng-abrecht-2015-better}, MoverScore~\cite{zhao-etal-2019-moverscore}, and $\mathrm{S}^3$~\cite{peyrard-eckle-kohler-2017-principled}, due to issues of reproducibility. Details in~\Cref{sec:6-chapter-summeval-repro}.

For the LLM-as-judge metrics, we experiment with the following models: Llama 8b~\cite{grattafiori2024llama3herdmodels}, Mistral 7b~\cite{Jiang2023Mistral7}, and Prometheus 7b~\cite{Kim2023Prometheus}.\footnote{\texttt{meta-llama/Llama-3.1-8B-Instruct}, \texttt{mistralai/Mistral-7B-Instruct-v0.3}, \texttt{prometheus-eval/prometheus-7b-v2.0}} We choose these models as they are popular open-source models that have been shown to be effective in LLM-as-judge settings~\cite{Zheng2023JudgingLLM,Kim2023Prometheus}. We use a temperature of $0$ and report the prompts and rubrics used in~\Cref{sec:ch-7-prompts}. 



For each metric, we test all applicable settings: reference-based, reference-adjusted, and reference-free. The reference-based setting uses the gold references provided in SummEval, while the reference-adjusted setting uses machine-generated summaries, either abstractive or extractive, as approximate references. The intuition is that what other algorithms choose to include is a useful content signal, even without an exact stylistic match. For each evaluated model, we randomly sample $K$ references from the remaining models in the dataset. We evaluate $K = [1, 6, 11]$ for both reference-based and abstractive reference-adjusted settings, and $K = [1, 6, 8]$ for extractive reference-adjusted settings, as SummEval contains only 8 extractive models. Full list of models is in~\Cref{tab:summeval_models}.
\input{research-chapters/6-chapter/5-LLM-as-Judge}
\section{Conclusion}


In this work, we introduced reference-adjusted evaluation, a new paradigm for LLM-as-judge summarization evaluation that bridges the gap between the expensive reference-based setting and the weaker reference-free setting. Our experiments across three open-source LLMs showed that reference-adjusted evaluation improves over the reference-free setting and, in some cases, approaches or outperforms the gold reference-based setting, especially when only a single gold reference is available. These results suggest that researchers can meaningfully improve their evaluations on new test sets without incurring the full cost of human annotation by leveraging either machine generated summaries or naturally occurring data. Our reference-adjusted paradigm provides a practical and scalable path toward more robust automatic evaluation of summarization systems.


\section{Limitations}
Although our results demonstrate that reference-adjusted evaluations consistently outperform reference-free evaluations, several limitations should be considered. First, all experiments are conducted on a single dataset and language. It remains an open question whether our findings generalize to other domains or languages. Second, we experiment with relatively small open-source models, and the behavior of larger or proprietary models may differ. Given that stronger judges tend to perform better overall~\cite{Zheng2023JudgingLLM,Kim2023Prometheus},
the gains we report on small open-source models may represent a lower bound on what is achievable with larger judges.

\section{Ethical Considerations}
The primary ethical consideration in this work concerns the use of
LLMs as automated judges, which carries the risk of encoding and amplifying
the biases present in the underlying models, including biases related to
writing style, dialect, topic, or cultural framing. Reference-adjusted
evaluation inherits this risk and may introduce an additional source of bias
through the machine-generated references themselves. We mitigate this concern by only including synthetic references from models other than the one being evaluated. Researchers applying reference-adjusted evaluation in practice should carefully consider the provenance and potential biases of the related references they use, and should not treat automated evaluation scores as a substitute for human judgment in high-stakes domains.



